\chapter{Instruction Scheduling}\label{ch:instruction-scheduling}

Processors overlap work across cycles, so instruction order decides how fast a program runs.
Two orders can compute identical results with very different runtimes when one stalls waiting for operands and the other fills those gaps with independent work.
Instruction scheduling reorders operations to hide latencies, fill issue slots, and respect resource limits without changing meaning.
It is the back-end pass where the abstract notion of dependence meets the concrete timing of a pipeline.

\section{Dependences and Machine Models}\label{sec:ch26-dependences}

A schedule is legal when it preserves every dependence and resource constraint.
\term{Instruction scheduling} orders machine operations in time subject to those constraints.
\term{Data dependence} means one operation reads a value another writes, and covers read-after-write, write-after-read, and write-after-write relations.
\term{Latency} means the cycles before a result is usable, while \term{throughput} means how often an operation can issue.
\term{Issue width} counts how many operations can start per cycle, and \term{functional unit} names the hardware that executes a class of operations.
A \term{dependence DAG} has one node per instruction and one edge per dependence labeled with the required delay.

Consider three operations where the first loads a value, the second adds to it, and the third is independent.
Scheduling the add immediately after the load stalls for the load latency.
Moving the independent operation between them hides part of that stall.
No dependence is violated because the moved operation touches different registers.
The same idea scales to loops, where software pipelining overlaps iterations, and to wide-issue machines, where the scheduler packs independent operations into the same cycle.
The \code{load}, \code{add}, and \code{mul} names below are \code{machine opcodes} whose latencies the scheduler reads from the machine model.
\begin{irlisting}
; dependence DAG in textual form
n1: r1 = load [p]      ; latency 3
n2: r2 = add r1, r0    ; depends on n1 with delay 3
n3: r3 = mul r4, r5    ; independent of n1 and n2
\end{irlisting}
The naive order \code{n1, n2, n3} stalls two cycles between \code{n1} and \code{n2}.
The order \code{n1, n3, n2} issues the load at cycle $0$, the independent multiply at $1$, and the dependent add at $3$.
Cycle $2$ is still idle because the load-to-add latency is three cycles.
The result bits are identical.
Only the timing differs.

Scheduling scope determines how much freedom the pass has.
\term{Basic-block scheduling} reorders within one block and is always safe to attempt.
\term{Superblock} and \term{trace scheduling} work across biased branches by treating a hot path as one region with compensation code on the exits.
\term{Software pipelining} schedules loop bodies so that operations from different iterations overlap in a steady state.
Pre-pass scheduling runs before register allocation to expose parallelism, while post-pass scheduling runs after allocation when true machine latencies and spill code are known.
Pre-pass scheduling risks increasing register pressure; post-pass scheduling cannot rename to break false dependences.
Many compilers run both lightly rather than one aggressively.
Figure~\ref{fig:ch26-dag} shows the dependence DAG for the example above.
Table~\ref{tab:ch26-scheduler-scope} compares the scheduling scopes.

\begin{figure}[htbp]
\centering
\begin{tikzpicture}[node distance=1.2cm]
\node[basic block] (n1) {n1: load};
\node[basic block, below left=of n1] (n2) {n2: add};
\node[basic block, below right=of n1] (n3) {n3: mul};
\draw[cfg edge] (n1) -- (n2);
\end{tikzpicture}
\caption{The load feeds the add while the multiply stays independent, so the scheduler may place the multiply between them without violating any edge.}
\label{fig:ch26-dag}
\end{figure}

\begin{table}[htbp]
\caption{Scheduling scopes with their freedom and main risk.}
\label{tab:ch26-scheduler-scope}
\begin{tabular}{lll}
\toprule
Scope & Freedom & Main risk \\
\midrule
Basic block & Local reorder only & Small hiding window \\
Superblock and trace & Cross-branch motion & Compensation code \\
Software pipeline & Cross-iteration overlap & Code growth, pressure \\
Pre-allocation & Full renaming room & Pressure blowup \\
Post-allocation & Exact latencies known & False dependences stuck \\
\bottomrule
\end{tabular}
\end{table}

Resource modeling decides how faithfully the scheduler tracks the machine.
A simple model counts issue slots per cycle and functional-unit occupancy.
A precise model adds pipeline bypasses, asymmetric ports, cracked micro-operations, and memory-disambiguation stalls.
Reservation tables record which resources each operation needs in each cycle after issue.
Overly simple models schedule for a machine that does not exist; overly precise models become brittle across chip revisions.
Most compilers ship a mid-level model tuned against benchmarks and refine the worst outliers by hand.
\begin{ccnote}[Model the machine you ship]
A scheduler tuned to last year pipeline makes this year code slower when port mappings change.
Recalibrate latencies and port assignments per subtarget and keep a fallback generic model.
A wrong table pessimizes every hot loop.
\end{ccnote}

\section{List Scheduling}\label{sec:ch26-algorithm}

\term{List scheduling} is the standard heuristic because optimal scheduling is NP-hard for essentially all realistic models.
It simulates cycles forward (or backward), keeps a ready list of operations whose predecessors are scheduled, and picks the highest-priority ready operation for each slot.
Priorities usually combine the critical-path height, the number of successors, and resource-fit tie-breakers.
Forward scheduling prioritizes hiding consumer stalls, while backward scheduling prioritizes exposing producer flexibility; compilers pick per target or run both and keep the better result.
Algorithm~\ref{alg:ch26-list} states the forward form for a single region with a reservation-table resource check.

\begin{algorithm}[htbp]
\caption{Forward list scheduling for one region}\label{alg:ch26-list}
\begin{pseudocode}
\Require complete dependence DAG $D$, resource table $R$, edge delays, individually schedulable operations
\Ensure legal schedule $S$ for the region
\State compute priority $H(n)$ as the critical-path height of $n$
\State $cycle \gets 0$, $Ready \gets$ nodes with no predecessors
\While{unscheduled nodes remain}
\State release each unscheduled node once all predecessors are placed and their edge delays have elapsed
\State select the highest $H(n)$ in $Ready$ that fits $R$ at $cycle$
\If{such $n$ exists}
\State place $n$ at $cycle$ and reserve its resources in $R$
\State remove $n$ from $Ready$
\Else
\State advance $cycle$ by one idle step
\EndIf
\EndWhile
\State emit $S$ in cycle order with exact slot assignments
\end{pseudocode}
\end{algorithm}

The algorithm assumes the dependence DAG is complete: every read-after-write, write-after-read, and write-after-write edge is present with its delay, memory dependences are disambiguated conservatively, and control dependences bound the region.
It also assumes the resource table answers fit queries in near-constant time.
Let $N$ be operations, $E$ dependence edges, $Q$ candidate fit checks, $W$ cost per check, and $T$ cycle advances.
With event-based readiness updates and a priority queue, bookkeeping costs $O(N\log N+E)$, but fit search and cycle advancement add $O(QW+T)$ for this implementation.
Repeatedly scanning blocked ready operations can make $Q$ much larger than $N$; large latencies likewise increase $T$ unless the scheduler jumps to the next event.
Priority computation costs $O(N + E)$ by longest-path traversal.
Priority-queue insertion and removal cost $O(N\log N)$ when each operation enters and leaves once; searching past resource-blocked candidates is additional work.
In practice the resource check dominates, so implementations cache port masks and precompute operation classes.
% TODO(cite): need a primary reference for trace scheduling with compensation code.

Heuristic quality hinges on the priority function and the tie-breakers.
Critical-path height alone is often enough for in-order pipelines.
Wide out-of-order machines need additional bias toward memory operations, long-latency producers, and operations that relieve port pressure.
Register-pressure awareness adds another term that delays operations whose results extend many live ranges.
No single priority wins everywhere, so schedulers expose weights and tune them per subtarget against loop-heavy benchmarks.

\section{Worked Example}\label{sec:ch26-worked}

We schedule a five-operation block by hand on a single-issue machine with the stated latencies.
The operations and dependences are small enough to enumerate completely.
\begin{irlisting}
; latencies: load 3, mul 2, add 1, store 1
a: r1 = load [p]
b: r2 = mul r1, r0
c: r3 = add r4, r5
d: r6 = mul r3, r2
e: store [q], r6
\end{irlisting}
Dependences: \code{b} waits $3$ cycles after \code{a}; \code{d} waits $2$ after \code{b} and $1$ after \code{c}; \code{e} waits $2$ after \code{d}.
Including each operation's own latency, critical-path heights are $1$ for \code{e}, $3$ for \code{d}, $5$ for \code{b}, $8$ for \code{a}, and $4$ for \code{c}.
The chain through \code{a, b, d, e} takes $8$ cycles.
List scheduling picks \code{a} before \code{c} because height $8$ exceeds $4$.
Cycle 0 issues \code{a}.
Cycle 1 has ready set \code{c} only, since \code{b} still waits for the load, so it issues \code{c}.
Cycle 2 has an empty ready set and idles.
Cycle 3 releases \code{b} and issues it.
Cycle 4 idles because \code{d} waits for the multiply latency.
Cycle $5$ issues \code{d}, cycle $6$ idles, cycle $7$ issues \code{e}, and the block completes at time $8$.
No legal schedule finishes sooner because the \code{a, b, d, e} chain alone needs that long.

Compare against the naive program order \code{a, b, c, d, e}.
That order must wait until cycle $3$ to issue \code{b}, then issues \code{c} at cycle $4$, \code{d} at $5$, and \code{e} at $7$.
It also completes at time $8$; rearranging the idle slots produces no elapsed-time gain in this example.
Verify legality of the final order \code{a, c, b, d, e}: every consumer issues after its producer delay, memory order between the load and store is preserved by disambiguation, and no resource is double-booked on the single-issue model.
Dual issue could place \code{a} and \code{c} together at cycle $0$, but the same dependence chain still prevents completion before time $8$.
Figure~\ref{fig:visual-schedule-timeline} places both single-issue schedules on one time axis.
The filled cells mark issue events, not the entire time an operation spends in the pipeline.
This distinction prevents confusing an empty issue slot with an idle processor: a previously issued load or multiply can still be in flight.

\begin{figure}[htbp]
  \centering
  \begin{tikzpicture}[x=11mm,y=10mm]
    \foreach \t in {0,...,7} {
      \node[font=\small] at (\t+0.5,2.2) {\t};
      \draw[black!35] (\t,0) rectangle +(1,0.8);
      \draw[black!35] (\t,1) rectangle +(1,0.8);
    }
    \node[font=\sffamily\small,anchor=east] at (-0.2,2.2) {Issue cycle};
    \node[font=\sffamily\small,anchor=east] at (-0.2,1.4) {Program order};
    \node[font=\sffamily\small,anchor=east] at (-0.2,0.4) {List order};
    \foreach \t/\op in {0/a,3/b,4/c,5/d,7/e} {
      \filldraw[fill=black!12] (\t,1) rectangle +(1,0.8);
      \node at (\t+0.5,1.4) {$\op$};
    }
    \foreach \t/\op in {0/a,1/c,3/b,5/d,7/e} {
      \filldraw[fill=black!12] (\t,0) rectangle +(1,0.8);
      \node at (\t+0.5,0.4) {$\op$};
    }
    \draw[cfg edge] (0,-0.6) -- (8,-0.6);
    \node[font=\small,anchor=north] at (4,-0.7) {Both complete at time 8; empty cells are idle issue slots.};
  \end{tikzpicture}
  \caption{Moving the independent add $c$ changes where the machine idles but cannot shorten the load--multiply--multiply--store dependence chain.}
  \label{fig:visual-schedule-timeline}
\end{figure}


That sensitivity to the model is exactly why the scheduler reads latencies from data rather than hard-coding them.

\section{Implementation Notes}\label{sec:ch26-implementation}

Production schedulers split the work across three places: pre-allocation scheduling for parallelism, register allocation with pressure awareness, and post-allocation scheduling for exact stalls and bundling.
Very long instruction word and in-order targets lean on the scheduler heavily and add bundling, packetization, and explicit no-operation insertion after scheduling.
Out-of-order targets lean on it less because hardware reorders dynamically, but they still care about instruction-cache layout, micro-operation fusion, and port balance.
Loop scheduling uses modulo scheduling with a computed initiation interval, kernel-only code generation, and prologue and epilogue handling; when the loop resists pipelining, partial unrolling plus block scheduling is the fallback.

Memory dependences deserve the most testing.
Static alias analysis disambiguates most stack and global references, while unknown pairs stay ordered conservatively.
Hardware aliasing hints, non-temporal markers, and fence instructions add edges the DAG must carry explicitly.
A scheduler that moves a load above an unresolved store trades a rare stall for a wrong program.
Keep the disambiguation interface narrow so the scheduler cannot invent orderings the analysis never approved.

Compile-time cost stays proportional to region size times issue width for list scheduling, which is why regions are capped and huge unrolled bodies are scheduled in windows.
Profile hot loops to pick modulo candidates rather than pipelining every loop nest.
Measure with hardware counters, not just static cycle estimates, because cache misses and branch effects dwarf scheduler gains on many workloads.
Helpful organization of selector, scheduler, and allocator interaction appears in~\cite{morgan-1998-optimizing} and see~\cite{appel-1998-modern}.
Soundness discipline for motion across memory operations follows the validation view in~\cite{safonov-2010-trustworthy}.
Memory-ordering and sequencing rules that constrain legal motion are grounded in~\cite{wg14-2011-c11}.

\section{Exercises}\label{sec:ch26-exercises}

\begin{enumerate}
\item Schedule the block in Section~\ref{sec:ch26-worked} on a dual-issue machine. \label{ex:ch26-dual-issue} Show the cycle-by-cycle placement, count the saved cycles, and verify every dependence delay.
\item A load feeds an add with latency 3, and four independent operations are available. \label{ex:ch26-hide-latency} Construct the schedule that hides the stall fully and state how many independents are necessary in general.
\item A scheduler moves a load above a store with an unknown alias relation. \label{ex:ch26-memory-edge} Explain why the motion is illegal, add the missing edge, and reschedule the three-operation example.
\end{enumerate}

\section{Further Reading}\label{sec:ch26-reading}

Scheduling scope, software pipelining, and target interaction are developed in~\cite{morgan-1998-optimizing}.
A compact account of scheduler placement in the backend pipeline appears in~\cite{appel-1998-modern}.
Validation of motion across potentially excepting and memory operations is treated in~\cite{safonov-2010-trustworthy}.
The sequencing and memory-model constraints that limit reordering appear in~\cite{wg14-2011-c11}.
